\documentclass{beamer}%handout,pdftex,
\usetheme{Madrid} 
\graphicspath{ {./images/} }
\usepackage{hyperref}
\usepackage{bm}
\usepackage{caption}
\usepackage{subcaption}
\usepackage{amsmath}
\usepackage{tikz}
\usepackage{verbatim}
\usepackage{lipsum}
\usepackage{physics}


\definecolor{UCRed}{RGB}{224,1,34}
\usecolortheme[named=UCRed]{structure}
\useoutertheme{split}
%\useinnertheme{rectangles}
\setbeamertemplate{navigation symbols}{}

\setbeamercolor{block title}{bg=UCRed!30,fg=black}
\setbeamercolor{block body}{bg=UCRed!30,fg=black}

\bibliographystyle{plain}
\setbeamertemplate{bibliography item}[text]

\newcommand\blfootnote[1]{%
  \begingroup
  \renewcommand\thefootnote{}\footnote{#1}%
  \addtocounter{footnote}{-1}%
  \endgroup
}
\newif\ifplacelogo % create a new conditional
\placelogotrue % set it to true
\logo{\ifplacelogo\includegraphics[height=0.8cm]{UC_ID_PrimaryBlackRed.pdf}\fi} % replace with your own command




\title[$M(D^+_s) - M(D^+)$]{A more precise measurement of $m_{D^+_s} - m_{D^+}$ through $D^{+}_{(s)} \rightarrow K^{+}K^{-}\pi^{+}$ decay}
\author[John Bodenschatz]{John Bodenschatz}
\institute{Physics Capstone\\ \vspace*{1.2em} Advisor: Dr. Michael Sokoloff}
\date{April 26, 2021}
%\logo{\includegraphics[height=0.8cm]{UC_ID_PrimaryBlackRed.pdf}}

%\titlegraphic{
%\includegraphics[height=1.5cm]{lhcb_logo.jpg}\vspace*{13cm}
%}

%\addtobeamertemplate{navigation symbols}{}{%
%    \usebeamerfont{footline}%
%    \usebeamercolor[fg]{footline}%
%    \hspace{1em}%
%    \insertframenumber/\inserttotalframenumber
%}

\begin{document}

\frame{\titlepage
\begin{tikzpicture}[remember picture,overlay]
    \node[xshift=1.0cm,yshift=1.05cm] at (current page.south west){\includegraphics[width=.14\linewidth]{lhcb_logo.jpg}};
        \node[xshift=0.0cm,yshift=1.2cm] at (current page.south){\includegraphics[width=.12\linewidth]{nsflogo.png}};
\end{tikzpicture}

}


%\begin{frame}{Outline of talk}
%\tableofcontents
%\end{frame}




\section{Motivation}
%~~~~~~~~~~~~~~~~~~~~~~~~~~~~MOTIVATION~~~~~~~~~~~~~~~~~~~~~~~~~~
  
  
\subsection{QCD}
  \begin{frame}
    %\frametitle{QCD implications}
Why are meson mass difference measurements important to Quantum Chromodynamics (QCD)?
\vspace{1em}
\begin{block}{}
"...in a rather direct manner [meson mass differences] permit to extract the ratios of quark masses which are fundamental parameters of QCD.\cite{theory}"
\end{block}
\vspace{1em}
Meson mass differences can help determine fundamental constants such as $m_c/m_b$

{\large
\begin{center}
$\ket{D^+} = \ket{c\bar{d}}$ \\
$\ket{D^+_s} = \ket{c\bar{s}}$ 
\end{center}
}
  \end{frame}
  
  
  
  \subsection{Precision}
  \begin{frame}

    {\Large For the $m_{D^+_s} - m_{D^+}$ measurement:}
    \vspace{0.5em}
\begin{itemize}
\setlength\itemsep{1.1em}
{\large
\item Most recent precision measurement was conducted on LHC Run 1 data from 2011\cite{lhcb}. $98.68 \pm 0.03\text{(stat)} \pm 0.04\text{(syst)}$ MeV
\item Using data from Run 2 could reduce statistical error by a factor of 10+
}
\item Reduce systematic uncertainty 
\end{itemize}

\end{frame}


\section{Method}
%~~~~~~~~~~~~~~~~~~~~~~~~~METHOD~~~~~~~~~~~~~~~~~~~~~~~~~~~~~~~~~~

\subsection{Getting the data}

\begin{frame}
%\frametitle{Method: Getting data from CERN}

\begin{figure}
\includegraphics[width=1.0\linewidth]{lhcb_run_2_data_flow.png}
\end{figure}
%\vspace{0.5em}
{\Large \begin{center}
Run 2 (2017) $D^{+}_{(s)} \rightarrow K^{+}K^{-}\pi^{+}$ decay data: 1.28 TB


\blfootnote{\href{https://lhcb.github.io/starterkit-lessons/first-analysis-steps/run-2-data-flow.html}{Image from LHCb StarterKit Lessons}}
\end{center}}

\end{frame}


\subsection{Cuts}

\begin{frame}
%\frametitle{Method: Cuts}
\textbf{Cuts}: set of selection criteria the events must meet to be selected for fitting. \underline{Isolate the signal peaks from background noise.}


\vspace{1.5em}

Examples of cuts used:
\begin{itemize}
\item P($K^+$), P($K^-$) $\geq 10,000$ MeV
\item PID probability cut on $K^+$, $K^-$, $\pi^+$
\item $ D^{+}_{(s)} \rightarrow \phi(1020)\pi^+ \rightarrow K^{+}K^{-}\pi^{+}$
\end{itemize}

\vspace{1.5em}

After all cuts and variable reduction are applied: \textbf{380 MB} of data
\end{frame}




\begin{frame}
%\frametitle{Method: Dalitz Plot - $\phi(1080)$ Cut}
\begin{figure}
\includegraphics[width=0.77\linewidth]{dalitz.png}
\end{figure}
\end{frame}



%\placelogotrue



\section{Results}
%~~~~~~~~~~~~~~~~~~~~~~~~~RESULTS~~~~~~~~~~~~~~~~~~~~~~~~~~~~~~~~

\subsection{Mass difference measurement}

\placelogofalse
\begin{frame}
%\frametitle{Method: mass difference measurement}
\begin{figure}
\includegraphics[width=1.00\linewidth]{bigfit.png}
\end{figure}
\end{frame}


\begin{frame}
%\frametitle{Method: mass difference measurement}
\begin{figure}
\includegraphics[width=0.8\linewidth]{linfitonly.png}
\end{figure}
\end{frame}


\placelogotrue


\begin{frame}
\frametitle{Other $D^{+}_{(s)} \rightarrow K^{+}K^{-}\pi^{+}$ analyses}

\begin{table}
\bgroup
\def\arraystretch{1.3}
{\footnotesize
\begin{tabular}{l | c | c | c}

Organization &  Year & $m_{D^+_s} - m_{D^+}$ [MeV] & $D^+_s$ Signal Events\\
\hline \hline
BaBar\cite{babar} & 2002 & $98.4 \pm 0.1\text{(stat)}\pm 0.3\text{(syst)}$ & $47,794 \pm 311$\\ 
LHCb\cite{lhcb} & 2013 & $98.68 \pm 0.03\text{(stat)}\pm 0.04\text{(syst)}$ & $248,694 \pm 540$\\
LHCb - UC & 2020 & $98.97 \pm 0.1\text{(stat)}$ & $12,529 \pm 1602$\\ 
LHCb - UC & 2021 & $98.89 \pm 0.0036\text{(stat)} \pm 0.xxx\text{(syst)}$ & $8,524,400 \pm 17,505$\\ 

\end{tabular}
}
\egroup
\end{table}

\blfootnote{LHCb Preliminary}
\end{frame}




\begin{frame}
%\frametitle{Method: comparison picture}
\begin{figure}
\includegraphics[width=0.83\linewidth]{dp_mass_comparison.png}
\end{figure}
\end{frame}






\subsection{Future work}
  \begin{frame}
    %\frametitle{BLANK RESULTS FRAME}
\begin{itemize}
\setlength\itemsep{1.5em}
{\LARGE
\item Momentum scaling corrections
\item Calculate other systematic errors
\item Refine fit
\item Publish!
}
\end{itemize}
  \end{frame}


%linear sum
%1968.77 - 1869.80 = 98.97 ± 0.1 MeV/c^2
%quadratic sum
%1968.77 - 1869.80 = 98.97 ± 0.1 MeV/c^2
\section*{}
\subsection*{References}
\begin{frame}
\begin{thebibliography}{3}

\bibitem{theory}
Jose L. Goity, Chandana P. Jayalath,
\textit{Strong and Electromagnetic Mass Splittings in Heavy Mesons},
\href{https://arxiv.org/pdf/hep-ph/0701245.pdf}{\texttt{arXiv:hep-ph/0701245}}

\bibitem{lhcb} 
LHCb Collaboration, R. Aaij \textit{et al.},
\textit{Precision measurement of D meson mass differences},
\href{https://arxiv.org/pdf/1304.6865.pdf}{\texttt{arXiv:hep-ex/1304.6865}}

\bibitem{babar} 
BaBar Collaboration, B. Aubert \textit{et al.},
\textit{Measurement of $D^+_s$ and $D^{*+}_s$ production in B meson decays and from continuum $e^+e^-$ annihilation at $\sqrt{s} = 10.6$ GeV},
Phys. Rev. \textbf{D65} (2002) 091104, \href{https://arxiv.org/pdf/hep-ex/0201041.pdf}{\texttt{arXiv:hep-ex/0201041}}

\end{thebibliography}
\end{frame}

\end{document}